\hypertarget{appendix-4fbc03d20829}{}
\section*{Two's complement}
\section{Two's complement}\label{twos-complement}

\textbf{Curriculum depth:} Developed · \textbf{First seen:} C001 · \textbf{Latest development:} C002

\subsection{Definition and mechanism}\label{definition-and-mechanism}

An \(N\)-bit negative \(-x\) is residue \(2^N-x=\sim x+1\); the signed high-bit weight is \(-2^{N-1}\), giving one zero and an asymmetric range.

For raw unsigned value u in N bits, signed decoding is u when u\textless2ᴺ⁻¹, otherwise u−2ᴺ. Thus \texttt{0xAD} means 173−256=−83. Negation uses width-limited complement plus one modulo 2ᴺ. Signed interpretation and unsigned residue share bits but have different language arithmetic contracts.

\subsection{Prerequisites}\label{prerequisites}

\begin{itemize}
\tightlist
\item
  \hyperref[appendix-c99977d88b33]{Unsigned modular arithmetic} --- Modulo m retains the remainder class after division by m.
\item
  \hyperref[appendix-8f38927cfc54]{Bitwise NOT} --- Complement must specify width: eight-bit complement of \texttt{0x0F} is \texttt{0xF0}, while sixteen-bit complement is \texttt{0xFFF0}.
\end{itemize}

\subsection{Used by}\label{used-by}

\begin{itemize}
\item
  \hyperref[appendix-1d0dbe8924cc]{Fixed-point representation} --- uses this mechanism to define its representation, error or validation contract.
\item
  \hyperref[appendix-05f39b10364d]{Minimum signed value} --- uses two's complement as part of its representation or reasoning.
\item
  \hyperref[appendix-8d687c9c08ab]{Carry, overflow, negative, and zero flags} --- uses two's complement as part of its representation or reasoning.
\item
  \hyperref[appendix-1c7df340573e]{Subtraction through addition} --- uses two's complement as part of its representation or reasoning.
\item
  \hyperref[appendix-2ed044358b82]{Arithmetic policies} --- uses two's complement as part of its representation or reasoning.
\end{itemize}

\subsection{Recall and next depth}\label{recall-and-next-depth}

Explain two's complement without the label, work the example above, then state the assumption whose removal would change the result.

\textbf{Pending:} signed shifts, instruction semantics, overflow intrinsics, and wider arithmetic.

\subsection{Source and navigation}\label{source-and-navigation}

\href{Continuous\%20Notes/Computing\%20Foundations\%20-\%20From\%20Bits\%20to\%20C++\%20and\%20Systems\%20APIs\#From\%20a\%20signal\%20to\%20an\%20interpretation}{C001 teaching chapter} · \href{Continuous\%20Notes/Computing\%20Foundations\%20-\%20From\%20Bits\%20to\%20C++\%20and\%20Systems\%20APIs\#Arithmetic\%20becomes\%20a\%20policy}{C002 development} · \hyperref[appendix-fd090d982a15]{Part definitions and mechanisms} · \href{}{Dependency Map (02-Foundations-II, appendix-4ad47bdd269a)}
